\subsection{Probar con un archivo de entrada}
\label{alg:2-0}

Copiar el caso de ejemplo a mano cada vez que se corre la solución es lento y
propenso a errores. Lo práctico es dejarlo en un \texttt{input.txt} y que un script
aparte se lo entregue al programa por \texttt{stdin}. Lo importante es que la
solución NO se toca: lee de \texttt{stdin} igual que en el juez, así que el archivo
que se prueba es exactamente el que se envía.

\subsubsection{Lo mínimo: redirección de la terminal}

Sin script y sin nada que instalar, la terminal ya conecta un archivo a
\texttt{stdin} con \texttt{<}:

\begin{lstlisting}[language=bash]
python sol.py < input.txt
pypy3 sol.py < input.txt            # PyPy, varias veces mas rapido
g++ -O2 -std=c++17 -o sol sol.cpp && ./sol < input.txt
\end{lstlisting}

\vspace{0.3cm}
\subsubsection{Script de apoyo: un solo comando para Python y C++}

El script siguiente recibe el nombre de la solución, le pasa \texttt{input.txt},
imprime la salida y el tiempo, y si se le da un archivo con la salida esperada la
compara. Si la solución es \texttt{.cpp} la compila antes, de modo que el equipo usa
el mismo comando sin importar el lenguaje.

\begin{lstlisting}[language=bash]
python correr.py sol.py                    # usa input.txt
python correr.py sol.py caso.txt           # otra entrada
python correr.py sol.py caso.txt esp.txt   # compara con la esperada
python correr.py sol.cpp caso.txt esp.txt  # compila con g++ y ejecuta
\end{lstlisting}

\begin{lstlisting}[language=Python]
import subprocess, sys, os, time

sol = sys.argv[1]                        # sol.py o sol.cpp
ent = sys.argv[2] if len(sys.argv) > 2 else "input.txt"

if sol.endswith(".py"):
    cmd = [sys.executable, sol]
else:                                    # C++: compilar primero
    exe = os.path.abspath(sol.rsplit(".", 1)[0])
    if os.name == "nt": exe += ".exe"    # ruta absoluta: Windows no busca en .
    if subprocess.run(["g++", "-O2", "-std=c++17",
                       "-o", exe, sol]).returncode:
        sys.exit("error de compilacion")
    cmd = [exe]

t = time.perf_counter()
with open(ent, "rb") as f:               # el archivo ENTRA por stdin
    r = subprocess.run(cmd, stdin=f, capture_output=True)
dt = time.perf_counter() - t

sal = r.stdout.decode()
print(sal, end="" if sal.endswith("\n") else "\n")
if r.stderr: print("[stderr]", r.stderr.decode())
print("-- %.3fs  codigo %d" % (dt, r.returncode))

if len(sys.argv) > 3:                    # salida esperada
    esp = open(sys.argv[3]).read()
    print("OK" if sal.split() == esp.split()
          else "DIFERENTE, se esperaba:\n" + esp)
\end{lstlisting}

La comparación usa \texttt{.split()} en ambos lados, así que ignora espacios y
saltos de línea sobrantes, que es como corrige la mayoría de jueces. El tiempo
impreso sirve para saber si se está cerca del límite, y el código de salida delata
un error de ejecución aunque el programa haya alcanzado a imprimir algo.

Solo usa \texttt{subprocess}, \texttt{sys}, \texttt{os} y \texttt{time}, los cuatro
de la biblioteca estándar: no hay nada que instalar. Como la solución se ejecuta con
\texttt{sys.executable}, el intérprete que mide es el mismo con el que se lanza el
script; para cronometrar igual que el juez cuando se va a enviar en PyPy, hay que
correrlo con \texttt{pypy3 correr.py sol.py}, no con \texttt{python}.

\vspace{0.3cm}
\subsubsection{Si se ejecuta con el botón Run del editor}

Cuando no hay dónde escribir argumentos ni redirección, la alternativa es redirigir
\texttt{stdin} desde el propio código, con una condición que lo desactiva sola en el
juez (allí no existe \texttt{input.txt}), de modo que tampoco hay que borrar nada
antes de enviar:

\begin{lstlisting}[language=Python]
import sys, os
if os.path.exists("input.txt"):      # en el juez no existe -> stdin real
    sys.stdin = open("input.txt")
\end{lstlisting}

\graybold{Cuidado:} lo que NO se debe hacer es leer el archivo directamente
(\texttt{open("input.txt").read()}) en lugar de reasignar \texttt{sys.stdin}: eso
obliga a editar la solución antes de enviarla, y si se olvida, el juez da error de
ejecución porque ese archivo no existe.
